\subsection*{An exact Clifford lift and the ordinary closed-square obstruction (NS-47)}

This subsection constructs a representation of the existing complete Weil
form. It does not establish the requested uniform lower bound. The
identities and the obstruction below are unconditional, using the archived
v1.36 and v1.50 results. The remaining estimate is identified explicitly.
There is no new window, tail metric or numerical claim; G2 and RH remain
open. All Hilbert inner products are linear in the first slot.

\paragraph{The two signed channels.}
Let $H=L^2(\mathbb R,dx)$, $\mathcal C=C_c^\infty(\mathbb R)$,
$I_a=(-a,a)$ and $\lambda=e^a$. Use precisely the positive even
weight $h$, probability measure $\nu$, and kernel $\rho$ of
\eqref{eq:v150-weight}--\eqref{eq:v150-global-energy}. In particular,
$I_h=\int \cosh(x/2)h(x)\,dx=1/\sqrt3$ and
$d\nu=I_h^{-1}\cosh(x/2)h(x)\,dx$.
Write $\ell_n=\log n$ and $w_n=\Lambda(n)/\sqrt n$.
Define the Hilbert spaces
\[
 E_+=L^2(\mathbb R^2,dx\,dy)\oplus
                 \bigoplus_{n\ge2} L^2(\mathbb R,dx),\qquad
 E_-=L^2_0(\nu)\oplus\mathbb C,
\]
where $L^2_0(\nu)$ denotes the mean-zero subspace. Zero-weight prime
channels are harmless. For $f\in\mathcal C$, put $g=f/h$ and
\begin{align}
 (D_0f)(x,y)&=\sqrt{\frac{\rho(|x-y|)h(x)h(y)}2}
                         (g(x)-g(y)),\notag\\
 (D_nf)(x)&=\sqrt{w_nh(x)h(x+\ell_n)}
                         (g(x+\ell_n)-g(x)),\notag\\
 Df&=(D_0f,(D_nf)_{n\ge2}),\label{eq:ns47-edge-map}\\
 Nf&=\left(\sqrt{\frac23}\left(g-\int g\,d\nu\right),
                  \sqrt2\int\sinh(x/2)f(x)\,dx\right).
 \label{eq:ns47-negative-map}
\end{align}
These maps are linear; the centering is taken on the whole probability
space, including outside the support of $f$.

\begin{proposition}[Exact signed lift, with a closed positive edge operator]
\label{prop:ns47-clifford-lift}
For the complete support-consistent form $q$ and every $f\in\mathcal C$,
\begin{equation}\label{eq:ns47-signed-lift}
 q[f]=\norm{Df}_{E_+}^2-\norm{Nf}_{E_-}^2.
\end{equation}
The maximal operator $D:H\supset\mathcal D(D)\to E_+$ given by
\eqref{eq:ns47-edge-map}, with domain where the squared-edge sum is
finite, is densely defined and closed. Consequently the graded operator
\begin{equation}\label{eq:ns47-dirac}
 \mathscr D=\begin{pmatrix}0&D^*\\D&0\end{pmatrix},\qquad
 \mathcal D(\mathscr D)=\mathcal D(D)\oplus\mathcal D(D^*)
 \quad\hbox{on }H\oplus E_+,
\end{equation}
is self-adjoint, anticommutes with $\Gamma=\operatorname{diag}(I,-I)$,
and has square $\operatorname{diag}(D^*D,DD^*)$ on its square domain.
Here $D^*D$ represents the positive edge energy, not the complete Weil
operator.

For an explicit Clifford module put $K=E_+\oplus E_-$, let
$i_\pm:E_\pm\to K$ be the canonical inclusions, and set
\[
 S f=(i_+Df,i_-Nf)\in K\oplus K,\qquad
 e_+=\begin{pmatrix}I&0\\0&-I\end{pmatrix},\quad
 e_-=\begin{pmatrix}0&I\\-I&0\end{pmatrix}.
\]
Then $e_+^2=I$, $e_-^2=-I$, $e_+e_-+e_-e_+=0$, and
\begin{equation}\label{eq:ns47-clifford}
 q[f]=\langle e_+Sf,Sf\rangle.
\end{equation}
Thus this is an explicit $\mathrm{Cl}_{1,1}$ representation with
infinite-dimensional coefficient space, not a representation of all
arithmetic channels by five scalar coordinates.
\end{proposition}
\begin{proof}
Proposition~\ref{prop:v150-critical-energy} gives finite separate energies
on $\mathcal C$ and exactly \eqref{eq:ns47-signed-lift}, including the
factor $1/2$ on the ordered continuous edges and the coefficient $2/3$
on the variance. To prove maximal closedness, suppose $f_j\to f$ in
$H$ and $Df_j\to F$ in $E_+$. Pass to a subsequence with almost
everywhere convergence of $f_j$. Pass further to a subsequence giving
almost everywhere convergence in every component of $Df_j$ (there are
only countably many). Translated null sets remain null, and the union
of the two coordinate null sets is null for product measure. Hence the
component limits equal the displayed differences of $f/h$ almost
everywhere. The diagonal $x=y$ is irrelevant. Thus $Df=F\in E_+$.
The domain contains $\mathcal C$ and is dense. The block-adjoint
calculation, using $D^{**}=D$, proves self-adjointness of
\eqref{eq:ns47-dirac}; see also the classical general construction in
\cite[Appendix A, Proposition A.1]{Holzmann2021NS47}. The grading and
Clifford relations follow by multiplication. No positivity of the
indefinite pairing in \eqref{eq:ns47-clifford} has been inferred.
\end{proof}

For each fixed $a$, zero extension identifies the maximal edge domain
with $\mathcal D(q_a)$, not merely with an unspecified closed extension
of the core. Indeed, on $[-a-1,a+1]$ the functions $h$ and $1/h$ are
bounded and Lipschitz, and $h$ has a positive minimum. The continuous
edges with $|x-y|<1$ therefore give, up to fixed-window constants and a
bounded ordinary-norm term, the killed logarithmic seminorm
\[
 \int_{|t|<1}\frac{\norm{f(\cdot+t)-f}^2}{|t|}\,dt.
\]
Multiplication by $h$ or $1/h$ preserves this domain: expand the
product difference, and bound the coefficient difference by a constant
times $|t|$. The Fourier multiplier of this seminorm is comparable,
after adding the ordinary norm, to $\log(2+|\xi|)$.
The continuous edges with $|x-y|\ge1$ and the entire prime-edge sum
are bounded ordinary-norm forms at fixed support. For the latter,
the Gaussian decay makes $\sum_n w_nh(x\pm\ell_n)/h(x)$ uniformly
finite for $x\in I_a$. The former follows from the integrable
off-diagonal kernel and the same decay. These comparisons prove
maximal-domain equality with the physical logarithmic form domain;
all their constants may depend on $a$. The operator $N_a$ obtained
by restricting $N$ is bounded in $L^2(I_a)$, since $h^{-1}$ is bounded
there. For example
\begin{equation}\label{eq:ns47-local-cost}
 \norm{N_af}^2\le C_a\norm f^2,\qquad
 C_a=\frac{2}{\sqrt3}\sup_{x\in I_a}\frac{\cosh(x/2)}{h(x)}
          +2\norm{\sinh(x/2)}_{L^2(I_a)}^2.
\end{equation}
This particular bound diverges with $a$; it is not the required uniform
floor. All exterior continuous edges and all exterior prime edges are
still in $D$: their energy is exactly
$\int_{I_a}(T_a/h)|f|^2$ in \eqref{eq:v150-edge-split}.
Only the original correlations have the strict cutoff $n<e^{2a}$.

For even $f$ the second component of $Nf$ vanishes; for odd $f$ its
first component is $\sqrt{2/3}f/h$ and its second must be retained.
The actual even source condition remains $\langle f,u_a\rangle=0$,
or $\int hg\overline{u_a}=0$; it is not the centering condition in
$N$. Parity orthogonality gives the two original blocks and their
sum without a missing cross term. A bounded invertible change of
Clifford coordinates satisfying $U^*e_+U=e_+$ preserves the displayed
form and cannot itself establish its sign.

\begin{proposition}[The exact contraction demand and its critical constant]
\label{prop:ns47-contraction}
Fix $C\ge0$, let $M_Cf=(Df,\sqrt C f)$, and let
$\mathcal H_C=\overline{M_C\mathcal C}\subset E_+\oplus H$.
The following are equivalent:
\begin{enumerate}[label=(\roman*)]
\item $q[f]\ge-C\norm f^2$ for every $f\in\mathcal C$;
\item there is a contraction $K_C:\mathcal H_C\to E_-$ such that
      $Nf=K_CM_Cf$ for every $f\in\mathcal C$.
\end{enumerate}
They are also equivalent to a common full-space lower floor on a
cofinal family of complete windows. If they hold, then
\begin{equation}\label{eq:ns47-positive-factor}
 q[f]+C\norm f^2=\norm{A_Cf}^2,\qquad
 A_C=(I-K_C^*K_C)^{1/2}M_C\big|_{\mathcal C}.
\end{equation}
The same range construction on each admissible window domain gives
the equivalent source-complement version of CAE. It requires one
common $C$ in both parities. No $K_C$ or that uniform estimate is
constructed here.

For this Gaussian edge map, any such full-space contraction must have
norm exactly one. The negative channel itself satisfies, on either
admissible parity space,
\begin{equation}\label{eq:ns47-negative-growth}
 \norm{N_a|_{\rm adm}}^2\ge
       [\kappa\exp(\gamma e^{2a})-B_a]_+.
\end{equation}
Thus merely bounding this channel in ordinary norm cannot give a
uniform error. More generally, a windowwise estimate
\begin{equation}\label{eq:ns47-strict-contraction}
 \norm{Nf}^2\le(1-\delta_a)\norm{Df}^2+C\norm f^2,
 \qquad 0\le\delta_a\le1,
\end{equation}
on either the even actual source complement or the odd sector forces
\begin{equation}\label{eq:ns47-critical-rate}
 \delta_a\le \frac{C+B_a}{\kappa}\exp(-\gamma e^{2a})
\end{equation}
for large $a$, with precisely $B_a,\kappa,\gamma$ of
Proposition~\ref{prop:v150-capacity-slack}.
\end{proposition}
\begin{proof}
The signed identity changes (i) into $\norm{Nf}\le\norm{M_Cf}$.
Define $K_C(M_Cf)=Nf$. This is well-defined, since the inequality
annihilates $N$ on the kernel of $M_C$, and is contractive; it extends
uniquely to $\mathcal H_C$. The converse and
\eqref{eq:ns47-positive-factor} follow by taking norms. This is the
elementary range-factorization argument underlying Douglas's
theorem~\cite{Douglas}; it does not assume a bounded global
operator $N$. Support consistency and the fixed-window form cores
give the asserted cofinal equivalence. The same proof works on each
source-admissible linear domain, without identifying these varying
domains as one global subspace. The source residual and
Corollary~\ref{cor:v136-complement-floor} then give the already known
conditional full-floor implication.

Use the unit vectors constructed in the proof of
Proposition~\ref{prop:v150-capacity-slack}. They are admissible in
each indicated parity sector and satisfy
$\norm{Df}^2\ge\kappa\exp(\gamma e^{2a})$ and $q_a[f]\le B_a$.
The identity $\norm{Nf}^2=\norm{Df}^2-q_a[f]$ proves
\eqref{eq:ns47-negative-growth}.
Equation~\eqref{eq:ns47-strict-contraction} would give
$q_a[f]\ge\delta_a\norm{Df}^2-C$, proving
\eqref{eq:ns47-critical-rate}. If $\norm{K_C}\le\theta<1$,
then $\norm{Nf}^2\le\theta^2(\norm{Df}^2+C\norm f^2)$ gives
the same contradiction with $\delta=1-\theta^2>0$ and ordinary
error $\theta^2C$. All support and source cross terms in these tests
are already controlled in v1.50. This reuses that analytic result;
it is not a new numerical experiment.
\end{proof}

\paragraph{A restriction on a global positive-square construction.}
The next obstruction concerns a single factor on the ordinary space
$H$. It does not concern the fixed-window self-adjoint operators,
an incompatible family of local factors, or a completion in a Weil
metric. ``Bounded remainder'' means a two-sided ordinary-norm bound
on a Hermitian quadratic form, including the special case of a
nonnegative error bounded by $C\norm f^2$.

\begin{lemma}[An unbounded positive edge at fixed support]
\label{lem:ns47-high-frequency}
In either parity there are unit vectors in a single
$C_c^\infty(I_b)$, for fixed $b>0$, on which $q$ tends to $+\infty$.
More precisely, for a fixed nonzero even $\varphi\in C_c^\infty(I_b)$,
the normalized functions $\varphi(x)\cos(Rx)$ and
$\varphi(x)\sin(Rx)$ have form $\log R+O_b(1)$ as $R\to\infty$.
\end{lemma}
\begin{proof}
At this fixed support the exact multiplier of
Proposition~\ref{prop:v135-both-parities} satisfies
$\beta_b(\xi)=\log(2+|\xi|)+O_b(1)$. The finite prime sum and
finite continuum integral are bounded, and the digamma asymptotic
gives this estimate uniformly on the real axis after enlarging the
constant on a compact interval. For a single modulation, translate
$\widehat\varphi$ by $R$ and split the integral at
$|\xi-R|=R/2$. Schwartz decay makes the exterior error $o(1)$,
while on the interior the logarithm is $\log R+O(1)$.
The modulations by $R$ and $-R$ have a mixed weighted Fourier
integral $o(1)$: their Schwartz peaks are separated by $2R$, and
the multiplier grows only logarithmically. The squared ordinary
norms of the cosine and sine tests tend to
$\norm\varphi^2/2$, with errors decreasing faster than any power
of $R$. Normalizing proves the assertion in each parity. These are
analytic tests, not additional computed windows.
\end{proof}

\begin{proposition}[No ordinary closable square with bounded remainder]
\label{prop:ns47-no-closed-square}
Use the complete form $q$ on $\mathcal C$ and the dense cutoff-radical
property of Lemma~\ref{lem:v136-total-radicals}. There do not exist
a Hilbert space $E$, a linear map $T:\mathcal C\to E$ closable
in ordinary $L^2(\mathbb R)$, and a bounded self-adjoint operator
$B$ on that space such that
\begin{equation}\label{eq:ns47-forbidden-square}
 q[f]=\norm{Tf}^2+\langle Bf,f\rangle
 \quad(f\in\mathcal C).
\end{equation}
The same assertion holds separately on either complete parity core.
More generally, if $T$ is ordinary closable and only the comparison
\begin{equation}\label{eq:ns47-lower-comparison}
 q[f]\ge\norm{Tf}^2-C\norm f^2\qquad(f\in\mathcal C)
\end{equation}
holds for a fixed $C\ge0$, then $\overline T$ is everywhere defined
and bounded, with $\norm{\overline T}\le\sqrt C$. In particular an
unbounded ordinary closable Dirac factor cannot even be a global
lower comparison with uniform ordinary error.
In particular $q$ has no semibounded closable extension in the ordinary
global $L^2$ space. If the contraction in
Proposition~\ref{prop:ns47-contraction} exists, the factor $A_C$ in
\eqref{eq:ns47-positive-factor} is necessarily nonclosable there.
No negative physical Weil direction is asserted.
\end{proposition}
\begin{proof}
Let $\mathcal R$ be the dense translate span of
Lemma~\ref{lem:v136-total-radicals}. For any fixed $r\in\mathcal R$,
its compact cutoffs $r_a$ satisfy $r_a\to r$ in $H$ and
$q[r_a]\to0$, the latter by their complete operator residual.
If \eqref{eq:ns47-forbidden-square} held, then
\[
 \norm{Tr_a}^2=q[r_a]-\langle Br_a,r_a\rangle
                   \longrightarrow-\langle Br,r\rangle.
\]
This is bounded. Choose a weakly convergent subsequence
$Tr_a\rightharpoonup v$. The graph of the closed operator
$\overline T$ is a closed linear subspace of $H\oplus E$ and
hence is weakly closed. It contains the weak limit $(r,v)$.
Consequently
\[
 r\in\mathcal D(\overline T),\qquad
 \norm{\overline T r}^2\le-\langle Br,r\rangle
                      \le\norm B\norm r^2.
\]
This bound holds for every $r$ in the dense linear space
$\mathcal R$. Applying it to differences, and using closedness
once more, extends $\overline T$ to all of $H$ with
$\norm{\overline T}\le\sqrt{\norm B}$.
Equation~\eqref{eq:ns47-forbidden-square} would therefore give
$q[f]\le2\norm B\norm f^2$ on the compact core, contradicting
Lemma~\ref{lem:ns47-high-frequency}.
For \eqref{eq:ns47-lower-comparison}, the same cutoff and weak-graph
argument gives $\norm{\overline T r}^2\le C\norm r^2$ on
$\mathcal R$, and the same closedness argument extends this bound
to all of $H$. This step does not need an upper bound on $q$.

For each parity use its dense reflected translate span from v1.36,
its same-parity cutoffs, and the corresponding high-frequency tests.
A semibounded closable form extension would have a closure whose
nonnegative constant shift is the square norm of a closed operator;
its restriction to $\mathcal C$ would yield
\eqref{eq:ns47-forbidden-square} with $B$ a scalar operator, which
has just been excluded. Finally, \eqref{eq:ns47-positive-factor}
has exactly that form with $B=-CI$, so its factor cannot be
closable. No assumption about a sign or about RH entered the proof.
\end{proof}

The two-sided remainder hypothesis matters for excluding every exact
representation \eqref{eq:ns47-forbidden-square}. A lower bound on an
additive signed remainder, or an upper bound on a subtracted cost,
instead gives only \eqref{eq:ns47-lower-comparison}, which forces the positive factor
to be bounded; it does not exclude all bounded comparisons. The
proposition forbids promoting the global geometric lift to an ordinary
closable positive square plus a bounded error. The dense-radical input
is inherited from v1.36;
this is a scoped consequence, not a claim of a new general theory of
Clifford algebras or of worldwide novelty. Its scope is distinct from
the vanishing-error positive-comparison statement of v1.36.

\paragraph{Named open input and next admissible construction.}
Call the unproved demand \emph{critical Clifford domination}:
construct, on the complete admissible domains, contractions
$K_{a,C}$ satisfying $N_af=K_{a,C}(D_af,\sqrt C f)$ with a common
finite $C$, in the even actual source complement and the odd sector.
This is exactly CAE in the displayed coordinates, not a new independent
hypothesis or progress on its truth. A successful use of these
coordinates must allow the cofinal contraction norms to approach one:
no window-independent strict contraction factor below one is available.
All exterior channels, the odd pole, and the actual source condition
must be retained. For a full global
factorization the positive factor must also accommodate ordinary
nonclosability. Window-dependent factors or a different completion
are not excluded, but would need an explicit ordinary-norm transfer
and the uniform estimate; their existence alone supplies neither.

What had to change is the proposal to obtain the full bound from one
ordinary closed Dirac square with bounded remainder. That construction
is excluded by Proposition~\ref{prop:ns47-no-closed-square}. The signed
Clifford lift exists exactly, and its positive edge block is closed,
but the required cancellation estimate has not been proved. This is
a failure of the specified construction, not a disproof of the
uniform floor or of RH. The broader geometric approach remains open;
G2 and RH remain open.
